
\documentclass[10pt,a4]{article}
\usepackage[english]{babel}
\usepackage[utf8x]{inputenc}
\usepackage{amsmath}
\usepackage{graphicx}
\usepackage[colorinlistoftodos]{todonotes}
\usepackage{hyperref}
\usepackage{geometry}
\geometry{top=3cm,left=2cm,right=2cm,bottom=3cm}
\usepackage[scaled]{helvet}
\usepackage[T1]{fontenc}
\renewcommand\familydefault{\sfdefault}


\author{Marco Ferraresi \and Edoardo Figini}
\date{\today}
\title{Rust Armor Separation Kernel » Port to RISC-V SMP}



\begin{document}
\maketitle
\tableofcontents

\section{Project data}

\begin{itemize}
\item
  Project supervisor: Prof. Vittorio Zaccaria

\item
Group composition:
\begin{center}
\begin{tabular}{lll}
Last and first name & Person code & Email address\\
\hline
  Ferraresi Marco & 10764701 & marco2.ferraresi@mail.polimi.it \\
  Figini Edoardo & 10771695 & edoardo.figini@mail.polimi.it \\
\end{tabular}
\end{center}

\item
Development tasks were \textbf{not} subdivided: the two of us worked on the
whole project together, in shared pair-programming sessions, from the first
``hello world'' on the board up to the final test suite. 

\item Links to the project source code:

\begin{itemize}
\item Repository: \href{https://github.com/polimi-aos-lab/aSK/tree/milkv-megrez}
{https://github.com/polimi-aos-lab/aSK}
\item Branch of this project: \href{https://github.com/polimi-aos-lab/aSK/tree/milkv-megrez}
{\texttt{milkv-megrez}}
\end{itemize}

\end{itemize}


\section{Project description}

The goal of this project is to port aSK to the Milk-V Megrez board, a new architecture 
featuring the RISC-V Hypervisor Extension and Memory Management Unit (MMU).

aSK (\emph{Armor Separation Kernel}) is a small type-1 hypervisor written in
Rust. It is a \emph{separation kernel}:
its only job is to run several execution environments on one chip and guarantee
that they cannot interfere with each other.

Our project was to make it run on a Milk-V Megrez board, based on
the ESWIN EIC7700X SoC (four harts, ISA \texttt{rv64imafdch}) and
to make it use the \textbf{native virtualization support} of RISC-V instead of
the hand-made approach: the hypervisor runs in HS-mode on top of standard
OpenSBI firmware, the guests run in VS/VU-mode, and isolation is enforced by
G-stage page tables, the second address translation stage that only the
hypervisor controls.

We faced two main challenges on this project:
\begin{itemize}
\item The \textbf{privilege level}: from M-mode to HS-mode, which changes the
control registers, the instruction used to enter a guest, and the way timers
and inter-processor interrupts are obtained.
\item The \textbf{memory model}: from PMP regions to page tables with
two translation stages.
\end{itemize}

This project is important for the AOS course because we got to apply the concepts about
virtualization and memory protection that we learned in class to a real board and kernel.

\subsection{Design and implementation}

aSK is organised in three layers: an architecture-agnostic \emph{kernel}, a
\emph{VHAL} layer of traits that states what an architecture must provide, and
the \emph{arch} implementations. Porting aSK therefore means implementing about
ten traits (\texttt{CpuOps}, \texttt{VmmOps}, \texttt{MpuOps},
\texttt{InterruptControllerOps}, ...) without touching the kernel.
All our code lives in a new module,
\texttt{sk/src/arch/riscv64h/}, plus the board crate
\texttt{boards/milkv-megrez-bxxx/}.

The module is selected by a \textbf{Cargo feature} (\texttt{hsmode} against
\texttt{mmode}) rather than by a custom target, since Spike and the Megrez share
the same target flag.

\subsubsection*{Memory model}

The first challenge we faced is how to make the region-based memory API survive on
top of page tables. 

The kernel talks to an MPU-shaped interface, so each method
was reimplemented in terms of translation: hypervisor regions become entries in
an Sv39 identity map rooted in \texttt{satp}, VM regions become entries in the
per-VM Sv39x4 G-stage rooted in \texttt{hgatp}, and enabling protection means
switching \texttt{satp} from Bare to Sv39. 

The hypervisor identity map covers the first 4~GiB with four 1~GiB superpages, 
which keeps the page-table footprint to a single table; page tables come from static pools in the BSS. 
Each VM gets its own G-stage, restricted to its own RAM window and tagged with its own VMID.

\subsubsection*{Virtual Machine Monitor}

From a guest's point of view aSK is the firmware: every \texttt{ecall} from VS-mode 
traps into HS-mode and is served by
our SBI implementation, with timer requests forwarded to the real OpenSBI underneath.
Devices a guest does not own are given to it by \emph{trap-and-emulate}: the
MMIO address is left out of the G-stage, the resulting guest-page fault is
decoded from \texttt{htval} and \texttt{htinst}, and an NS16550 UART is emulated
in software.

Since \texttt{htinst} is not reliable in RISC-V H-extension V0.6, the hypervisor has to decode 
the instruction from the guest PC. A feature flag determines when the hypervisor runs on 
hardware that implements V0.6 (Section \ref{sec:hext}). 

\subsubsection*{Multi-hart boot}

Multi-hart bring-up follows the OpenSBI model, where only one hart starts and
the others must be woken through the HSM extension. This required a second
assembly entry point for secondary harts that skips the BSS clear (the master
has already filled global state by then) and a runtime election of the master
hart, because the boot hart is chosen by the firmware and is not predictable.

\subsubsection*{VM isolation}

Isolation is demonstrated by running two VMs, in disjoint RAM windows, on
different harts, each with its own G-stage and VMID, while one of them
deliberately reads the other's memory. The resulting fault is not treated as a
fatal error: it is logged and \textbf{only the offending VM is stopped}, which is
precisely the semantics that distinguishes a separation kernel from an ordinary
hypervisor.

Finally, VMs can own \textbf{real} devices: a VM declares an interrupt in its
static configuration, the kernel enables that PLIC source on the context of the
hart hosting it, and when the device fires, the interrupt is claimed in HS-mode,
matched against the routing table and injected into the owning guest as a
virtual external interrupt through \texttt{hvip}.

\subsubsection*{RISC-V H-extension V1.0 vs V0.6}
\label{sec:hext}

Milk-V Megrez implements the Hypervisor Extension in a pre-ratified version (V0.6).

The differences affect CSR behaviour and presence

\begin{table}[h!]
\centering
\begin{tabular}{ l p{4.5cm} p{7.0cm} }
\textbf{Feature} & \textbf{Issue} & \textbf{Solution} \\
\hline
\texttt{htinst} & Present but undocumented; cannot depend on it & 
Feature flag and instruction fetched from guest PC \\
\texttt{vstimecmp} & Not present & This was not an issue because timer are set using the SBI layer \\
\texttt{hstatus.VGEIN} & In V0.6 not implemented, reserved on megrez & This was not an issue because
we do not use guest external interrupts, we just needed to write 0 during boot \\
\texttt{htval} & Present but undocumented & No solution, we have to verify on the hardware if some 
issues arise \\
\texttt{henvcfg} & Not present & This was not an issue because we do not touch it \\
\texttt{hgeie} / \texttt{hgeip} & underspecified in V0.6 & This was not an issue because we do not 
touch them \\
\end{tabular}
\end{table}

\section{Project outcomes}

\subsection{Concrete outcomes}

The deliverable is the \texttt{milkv-megrez} branch of the aSK repository:

\begin{itemize}
\item A new architecture module, \texttt{sk/src/arch/riscv64h/}, implementing the
  full VHAL contract for HS-mode: boot and per-hart setup, CPU operations,
  memory management on Sv39 page tables, the VMM with its hypervisor
  vector table and SBI dispatcher, the PLIC driver and the interrupt controller with virtual
  interrupt injection and physical-IRQ routing.
\item A dual platform configuration in \texttt{sk/src/arch/riscv64h/boards/milkv\_megrez.rs}:
  the real EIC7700X one, and a QEMU one that is reproducible on any
  machine. The QEMU configuration is used in the test suite.

  \href{https://github.com/polimi-aos-lab/aSK/commit/a67417556fbbf4272e4b69bc4d37a015b91c15fe}
  {Commit a674175 "[milkv-megrez] EIC7700X platform config"}
\item The board crate \texttt{boards/milkv-megrez-bxxx/}, where (in \texttt{main.rs}) 
  the two VMs are declared statically with their memory,
  their emulated UART and the physical interrupt one of them owns.

  \href{https://github.com/polimi-aos-lab/aSK/commit/0de73bdfbcdef937a27c4903f2ae5304432a0d4c}
  {Commit 0de73bd "[milkv-megrez] multi-hart multi-VM boot flow"}
\item Integration into the repository test infrastructure: a Docker image with
  QEMU, OpenSBI and U-Boot, and two smoke-test scenarios:
  \begin{itemize}
  \item one covering the boot chain up to the point where the guests start.
  \item one covering the VM layer
  (MMIO emulation, interrupt injection, physical IRQ routing and end-of-interrupt,
  isolation violation and clean shutdown) 
  \end{itemize}
  
  run with \texttt{make test-riscv-h}.

  \href{https://github.com/polimi-aos-lab/aSK/commit/92be7d80fde63282afa43e8bfbbe2b73be3b400b}
  {Commit 92be7d8 "[milkv-megrez] integrated milkv-megrez in test suite"}
\end{itemize}

All eight phases of the roadmap agreed with the supervisor are complete and
verified on the emulated platform, and the first four (boot, console, CPU,
MMU with the identity map) were also validated on the physical board.

What remains is re-running the complete integrated flow on the hardware with the real
platform configuration; this is where the differences of the EIC7700X
implementation of the Hypervisor Extension, which predates the ratified
specification, are expected to show up. 

\subsection{Learning outcomes}

Tools we learned to use and knowledge we gained during this project:
\begin{itemize} 
  \item Rust programming language: 
  we both started this project with a surface level knowledge of Rust.
  This project gave us both the opportunity to explore the language (trait-implementation paradigm
  and safety oriented approach) and the ecosystem (Cargo, crates.io, Rustup, etc.) in depth.
  \item Privileged RISC-V:
  we were already familiar with the user-level RISC-V ISA from previous courses, this project 
  gave us the opportunity to explore how RISC-V works beyond the simple pipeline we had studied
  (control status registers, traps, interrupts, memory management).
  \item QEMU: 
  we learned how to configure QEMU to emulate a RISC-V board like Milk-V Megrez.
  \item Virtualization:
  we got to explore hands-on how hardware assisted virtualization enables virtual machines
  to run efficiently at near-native speed.

  We learned how the hypervisor does not actively control VMs, but handles exceptions raised
  by the hardware, allowing for simpler design and better performance.
\end{itemize}

We took this project as an opportunity to enrich our knowledge of low-level systems,
a subject we are both interested in and wanted to explore more.
Specifically, Marco was really interested 
in learning about Embedded Rust, since it is becoming a more and more popular language
with an active community and a growing ecosystem. Edoardo was more interested in developing a low-level
system and writing software that interacts with the hardware. Along with that, he wanted to
become more familiar with the RISC-V architecture as it is growing in popularity and it's exiting
the niche academic world into commercial applications.
Both of us will bring the knowledge we gained in this project into our future personal and
professional projects.

\subsection{Existing knowledge}

Some courses we followed have helped us doing this project:
\begin{itemize}
  \item Embedded Systems: theoretical concepts in line with AOS course, but with a focus on 
  firmware and hardware interaction: this helped us as this project was focused on bridging the gap
  between the hardware and the higher level software (the hypervisor).
  \item Sensor Systems and Design of Hardware Accelerators: in these courses we had to interact 
  with a board using only UART for debugging without fancy tools like \texttt{gdb} as we used in high level
  software development.
\end{itemize}

\subsection{Problems encountered}

\begin{itemize}
\item \textbf{Total silence at the first boot on the board.} The binary was
  loaded and started, and nothing happened: no output, no error. The cause was
  the firmware boot protocol: the U-Boot \texttt{go} command is a raw jump and
  does not set the registers that carry the hart identifier, so our code read
  garbage, concluded it was a secondary hart, and waited forever for a master
  that would never arrive. The solution came by looking at the U-Boot documentation online
  and, after a lengthy trial and error process, we concluded the correct command should have 
  been \texttt{bootelf}. This meant that U-Boot accepts elf binaries, so
  we didn't have to provide the raw binary.

\item \textbf{The boot hart is not predictable.} Our first multicore code
  assumed hart 0 would be the master. The firmware picks a different hart on
  different runs, which deadlocked the rendez-vous. We replaced the assumption
  with a runtime election, using an atomic compare-and-exchange on a variable
  placed in initialised data so that it survives the zeroing of the BSS.

\item \textbf{Unsynchronized console.} With several harts printing to a serial
  port with no locking, output lines from different cores are interleaved
  character by character. This is cosmetic when reading logs by eye, but it made
  automated test patterns unreliable. We synchronized the hypervisor's own output with a per-line
  lock, which also made the interesting events assertable in the test suite.
  Guest output is still interleaved, since the locking mechanism we implemented lives
  in the hypervisor and is not available to the guest.
\end{itemize}

None of us had encountered these specific problems before. What they have in
common is that the system does not show what is wrong: there is no error
message, no exception and often no output at all, so we could not rely
on a traditional debugging tools. We had to understand the behaviour of our code 
through rudimentary "printf" debugging. 

\section{Honor Pledge}

We pledge that this work was fully and wholly completed within the criteria
established for academic integrity by Politecnico di Milano (Code of Ethics and
Conduct) and represents our original production, unless otherwise cited.

We also understand that this project, if successfully graded,  will fulfill part B requirement of the
Advanced Operating System course and that it will be considered valid up until
the AOS exam of Sept. 2026.

\begin{flushright}
Marco Ferraresi, Edoardo Figini
\end{flushright}


\end{document}
